% Options for packages loaded elsewhere
\PassOptionsToPackage{unicode}{hyperref}
\PassOptionsToPackage{hyphens}{url}
\documentclass[
  12pt,
]{ctexart}
\usepackage{xcolor}
\usepackage{amsmath,amssymb}
\setcounter{secnumdepth}{-\maxdimen} % remove section numbering
\usepackage{iftex}
\ifPDFTeX
  \usepackage[T1]{fontenc}
  \usepackage[utf8]{inputenc}
  \usepackage{textcomp} % provide euro and other symbols
\else % if luatex or xetex
  \usepackage{unicode-math} % this also loads fontspec
  \defaultfontfeatures{Scale=MatchLowercase}
  \defaultfontfeatures[\rmfamily]{Ligatures=TeX,Scale=1}
\fi
\usepackage{lmodern}
\ifPDFTeX\else
  % xetex/luatex font selection
\fi
% Use upquote if available, for straight quotes in verbatim environments
\IfFileExists{upquote.sty}{\usepackage{upquote}}{}
\IfFileExists{microtype.sty}{% use microtype if available
  \usepackage[]{microtype}
  \UseMicrotypeSet[protrusion]{basicmath} % disable protrusion for tt fonts
}{}
\makeatletter
\@ifundefined{KOMAClassName}{% if non-KOMA class
  \IfFileExists{parskip.sty}{%
    \usepackage{parskip}
  }{% else
    \setlength{\parindent}{0pt}
    \setlength{\parskip}{6pt plus 2pt minus 1pt}}
}{% if KOMA class
  \KOMAoptions{parskip=half}}
\makeatother
\usepackage{longtable,booktabs,array}
\usepackage{calc} % for calculating minipage widths
% Correct order of tables after \paragraph or \subparagraph
\usepackage{etoolbox}
\makeatletter
\patchcmd\longtable{\par}{\if@noskipsec\mbox{}\fi\par}{}{}
\makeatother
% Allow footnotes in longtable head/foot
\IfFileExists{footnotehyper.sty}{\usepackage{footnotehyper}}{\usepackage{footnote}}
\makesavenoteenv{longtable}
\setlength{\emergencystretch}{3em} % prevent overfull lines
\providecommand{\tightlist}{%
  \setlength{\itemsep}{0pt}\setlength{\parskip}{0pt}}
\usepackage{bookmark}
\IfFileExists{xurl.sty}{\usepackage{xurl}}{} % add URL line breaks if available
\urlstyle{same}
\hypersetup{
  hidelinks,
  pdfcreator={LaTeX via pandoc}}

\author{SCX}
\date{2026}

\begin{document}

\section{SCX 研究进展摘要 ---
2026年6月28日}\label{scx-ux7814ux7a76ux8fdbux5c55ux6458ux8981-2026ux5e746ux670828ux65e5}

\begin{quote}
写给导师看的 5 分钟速览。需要展开的部分随时说。
\end{quote}

\begin{center}\rule{0.5\linewidth}{0.5pt}\end{center}

\subsection{一句话}\label{ux4e00ux53e5ux8bdd}

\textbf{SCX (State-Conditioned
eXpertise)}：专家可靠性不是全局常数，而是数据状态的条件函数。这个框架回答了三个问题------哪些数据是噪声（Thm
1）、方法在什么条件下会失效（Thm 2）、噪声和困难样本为什么区分不了（Thm
3）。

\begin{center}\rule{0.5\linewidth}{0.5pt}\end{center}

\subsection{做了什么（2026.5 →
6）}\label{ux505aux4e86ux4ec0ux4e482026.5-6}

\subsubsection{理论：3
个定理，全部有完整证明}\label{ux7406ux8bba3-ux4e2aux5b9aux7406ux5168ux90e8ux6709ux5b8cux6574ux8bc1ux660e}

\begin{longtable}[]{@{}
  >{\raggedright\arraybackslash}p{(\linewidth - 4\tabcolsep) * \real{0.2400}}
  >{\raggedright\arraybackslash}p{(\linewidth - 4\tabcolsep) * \real{0.3600}}
  >{\raggedright\arraybackslash}p{(\linewidth - 4\tabcolsep) * \real{0.4000}}@{}}
\toprule\noalign{}
\begin{minipage}[b]{\linewidth}\raggedright
定理
\end{minipage} & \begin{minipage}[b]{\linewidth}\raggedright
核心结果
\end{minipage} & \begin{minipage}[b]{\linewidth}\raggedright
证明方法
\end{minipage} \\
\midrule\noalign{}
\endhead
\bottomrule\noalign{}
\endlastfoot
\textbf{Thm 1} 多专家一致性噪声检测 & F1 ≥ 1 - exp(-2MΔ²)/η，M≥20
名专家时 F1≥0.95 & Hoeffding + Chernoff + 条件期望分解 \\
\textbf{Thm 2} 弱特征失效下界 & F1 ≤ baseline + √(2δ)，特征互信息 δ→0
时退化到随机 & Fano 不等式 + 耦合论证 + Pinsker \\
\textbf{Thm 3} 噪声与困难不可区分 &
构造性反例证明：没有任何算法能区分噪声与困难样本 & Gödel 式构造 \\
\end{longtable}

定理文件总计约 1,700 行数学证明，位于 \texttt{theory/theorems/}。

\subsubsection{代码：SCX 框架
v0.4.0}\label{ux4ee3ux7801scx-ux6846ux67b6-v0.4.0}

\begin{longtable}[]{@{}
  >{\raggedright\arraybackslash}p{(\linewidth - 2\tabcolsep) * \real{0.5455}}
  >{\raggedright\arraybackslash}p{(\linewidth - 2\tabcolsep) * \real{0.4545}}@{}}
\toprule\noalign{}
\begin{minipage}[b]{\linewidth}\raggedright
指标
\end{minipage} & \begin{minipage}[b]{\linewidth}\raggedright
\end{minipage} \\
\midrule\noalign{}
\endhead
\bottomrule\noalign{}
\endlastfoot
Python 包 & \texttt{src/scx/} --- 6 子包、44 模块、\textasciitilde10,800
行 \\
功能模块 &
状态发现、专家注册/路由/冲突检测、数据估值（噪声/冗余/可学习性）、采集/压缩策略 \\
测试 & \textbf{427 个单元测试，全部通过} \\
子项目 & \texttt{scx-health/} 医学数据估值（MedMNIST,
F1=0.93）、\texttt{drug-module/} 药物靶点计算管线 \\
\end{longtable}

\subsubsection{实验（CPU
可做的部分）}\label{ux5b9eux9a8ccpu-ux53efux505aux7684ux90e8ux5206}

\begin{itemize}
\tightlist
\item
  ML 通用基准（UCI tabular）
\item
  CIFAR-10/100 噪声检测验证
\item
  AlN v3 两层 MLIP 分析
\item
  合成 2D 验证实验
\end{itemize}

\begin{center}\rule{0.5\linewidth}{0.5pt}\end{center}

\subsection{论文规划：5 篇}\label{ux8bbaux6587ux89c4ux52125-ux7bc7}

\begin{verbatim}
Paper 2 (最接近投稿)     EGP Gauge Fixing          npj Comput. Mater.    已有草稿+8图, 1-2周可投
    ↓
Paper 3 (理论先占坑)     噪声检测理论 (Thm 1-3)     arXiv → TMLR          定理证明完毕
    ↓
Paper 1 (旗舰)          SCX 多领域验证             Nature Comput. Sci.   等 GPU
    ↓
Paper 4                 压缩保真理论               JMLR                  定理需完善
    ↓
Gatekeeper 论文         数据策展-探索权衡           Nature Comput. Sci.   规划中
    ↓
Paper 5                 跨领域深挖 (Sim/Health)     待定                  远期
\end{verbatim}

\textbf{发表策略}：Paper 2 先发（纯应用，不依赖 GPU）→ Paper 3 上 arXiv
占理论坑 → 后续应用论文引用理论保证。

\begin{center}\rule{0.5\linewidth}{0.5pt}\end{center}

\subsection{当前状态}\label{ux5f53ux524dux72b6ux6001}

\textbf{阻塞}：GPU 未到，以下实验无法进行： - AlN v3 去噪重训（Paper 1
核心实验） - CIFAR 完整训练 - MedMNIST 强 backbone 实验

\textbf{可以推进的}（无 GPU 依赖）： 1. 🔴 Paper 2 收尾投稿 2. 🟡 Paper
3 整理成 arXiv 预印本（定理已就绪） 3. 🟡 UCI tabular 实验 4. 🟢 公开
MLIP 数据下载分析

\begin{center}\rule{0.5\linewidth}{0.5pt}\end{center}

\subsection{需要和导师确认的}\label{ux9700ux8981ux548cux5bfcux5e08ux786eux8ba4ux7684}

\begin{enumerate}
\def\labelenumi{\arabic{enumi}.}
\tightlist
\item
  \textbf{论文顺序}：Paper 2（纯 EGP gauge fixing）先投 npj
  是否合理？还是应该等 Paper 1（SCX 框架+多领域）完成后以主论文形式发？
\item
  \textbf{学术优先权}：Thm 1-3 是否现在就要上 arXiv
  占坑？还是等实验结果补齐再一起发？
\item
  \textbf{GPU
  时间线}：组内显卡预计什么时候能到位？如果太久，是否有替代方案（租云
  GPU、借隔壁组）？
\item
  \textbf{合作}：药物靶点计算管线（\texttt{drug-module/}）需要药学院/生科院合作方数据验证，是否有推荐的合作者？
\end{enumerate}

\begin{center}\rule{0.5\linewidth}{0.5pt}\end{center}

\emph{完。任何时候需要看代码、证明、实验数据，直接说文件名即可。}

\end{document}
